\documentclass[11pt,a4paper]{article}
\usepackage[utf8]{inputenc}
\usepackage[T1]{fontenc}
\usepackage{amsfonts}
\usepackage[french]{babel}
\frenchbsetup{StandardLists=true}
\usepackage{enumitem}
\usepackage{amssymb}
\usepackage{amsmath}
\usepackage[left=2cm,right=2cm,top=2cm,bottom=2cm]{geometry}
\usepackage{multirow}
\usepackage[dvipsnames]{xcolor}

\usepackage[T1]{fontenc}
\usepackage{fouriernc}
\usepackage{graphicx}

\usepackage{setspace}
\singlespacing

\usepackage{tcolorbox}
\usepackage{multicol}
\usepackage{caption}
\usepackage{fancyhdr}
\pagestyle{fancy}
\renewcommand\headrulewidth{1pt}
\fancyhead[L]{Lycée Denis Diderot}
\fancyhead[R]{Classe de 3PM}
\setcounter{page}{1}\fancyfoot[C]{page \thepage}
\fancyfoot[R]{Année 2025-2026}
\fancyfoot[L]{Alexis RUIZ}
\usepackage{multirow,tabularx,booktabs}
\newcommand{\cadreblanc}[1]{%
\medskip\framebox[\linewidth][c]{%
\rule{0mm}{#1mm}}\par
\medskip}

\usepackage{awesomebox}
\setlength{\aweboxvskip}{0mm}
\usepackage{pifont}
\usepackage{chemmacros}
\chemsetup{modules=all}
\setlength{\parindent}{0pt}
\setlength{\headheight}{14pt}

\renewcommand{\arraystretch}{1.8}

\frenchbsetup{StandardLists=true}

\begin{document}

%------------------------------------------------------------
% TITRE
%------------------------------------------------------------
\begin{center}
  {\LARGE\textbf{TP : Le chou rouge comme indicateur de pH}}
\end{center}

\hrule
\vspace{3mm}

\noindent Nom : \dotfill \quad Prénom : \dotfill

\vspace{3mm}
\hrule
\vspace{5mm}

%------------------------------------------------------------
% PARTIE 1 -- PRÉPARATION DU JUS DE CHOU ROUGE
%------------------------------------------------------------

\textbf{\large Préparation du jus de chou rouge -- l'indicateur coloré}

\vspace{3mm}

\begin{itemize}
  \item Hacher le chou et le mettre dans un bécher en pyrex avec environ 50 mL d'eau.
  \item Chauffer le mélange jusqu'à ébullition.
  \item Eteindre et laisser refroidir quelques minutes.
  \item Filtrer en utilisant un entonnoir et du papier filtre. Recueillir la solution filtrée dans un erlen-meyer. On obtient un liquide bleu sombre qui servira d'indicateur coloré.
  \item Introduire le liquide obtenu dans un flacon compte-gouttes.
\end{itemize}

\begin{tcolorbox}[colback=orange!10!white, colframe=orange!80!black]
\textbf{Appeler le professeur.}
\end{tcolorbox}

\vspace{5mm}

%------------------------------------------------------------
% PARTIE 2 -- APPLICATIONS
%------------------------------------------------------------

\textbf{\large Applications}

\vspace{5mm}

%--- A. Echelle de teintes (expérience soude/acide) ---
\textbf{Echelle de teintes}

\vspace{3mm}

\begin{enumerate}[label=\textbf{\arabic*.}]

  \item Verser dans un tube à essais 1 cm d'acide chlorhydrique de concentration $10^{-1}$~mol/L. Ajouter 5 gouttes de l'indicateur coloré. Ajouter progressivement de la soude à la concentration $10^{-1}$~mol/L.

  \vspace{2mm}
  \textbf{Qu'observe-t-on ? Interpréter.}

  \cadreblanc{30}

  \item À présent, faire l'opération inverse, c'est-à-dire :
  \begin{itemize}
    \item Verser dans un tube à essai 1 cm de la soude à la concentration $10^{-1}$~mol/L.
    \item Ajouter 5 gouttes de l'indicateur coloré.
    \item Ajouter progressivement d'acide chlorhydrique de concentration $10^{-1}$~mol/L.
  \end{itemize}

  \vspace{2mm}
  \textbf{Qu'observe-t-on ? Interpréter.}

  \cadreblanc{30}

\end{enumerate}

\vspace{3mm}

%--- B. Echelle des teintes (6 tubes) ---
\textbf{Echelle des teintes :}

\vspace{3mm}

Préparer 6 tubes à essais, et les disposer de la gauche vers la droite :

\vspace{2mm}

Verser dans chaque tube à essai, sur une hauteur d'environ 4 cm, la solution de pH correspondant.

\vspace{3mm}

\begin{center}
\begin{tabular}{|c|l|c|}
\hline
\textbf{Tube} & \textbf{Solution} & \textbf{pH} \\
\hline
1 & Acide chlorhydrique & 1 \\
\hline
2 & Vinaigre & 3 \\
\hline
3 & Solution tampon & 4 \\
\hline
4 & Eau d'Évian & 7 \\
\hline
5 & Solution tampon & 10 \\
\hline
6 & Hydroxyde de sodium & 13 \\
\hline
\end{tabular}
\end{center}

\vspace{3mm}

Verser environ 12 gouttes d'indicateur coloré préparé dans chaque tube. Compléter le tableau :

\vspace{3mm}

\begin{center}
\begin{tabular}{|c|c|c|c|c|c|c|}
\hline
\textbf{pH} & \textbf{1} & \textbf{3} & \textbf{4} & \textbf{7} & \textbf{10} & \textbf{13} \\
\hline
\textbf{Couleur} & & & & & & \\[14pt]
\hline
\end{tabular}
\end{center}

\vspace{3mm}

\begin{tcolorbox}[colback=orange!10!white, colframe=orange!80!black]
\textbf{Appeler le professeur.}
\end{tcolorbox}

\vspace{5mm}

%--- C. Evaluation du pH ---
\textbf{Evaluation du pH}

\vspace{3mm}

Verser dans un tube à essais 1 cm de la solution de pH inconnu. Ajouter quelques gouttes de l'indicateur coloré. Observer la couleur et encadrer le pH de la solution.

\vspace{5mm}

\begin{center}
\framebox[3cm]{\rule{0pt}{8mm}} \quad $<$ \quad pH \quad $<$ \quad \framebox[3cm]{\rule{0pt}{8mm}}
\end{center}

\vspace{7mm}

%------------------------------------------------------------
% PARTIE 3 -- REMISE EN ÉTAT
%------------------------------------------------------------

\textbf{\large 3. Remise en état du poste de travail}

\vspace{3mm}

\begin{itemize}
  \item Jeter la solution de chou rouge dans l'évier.
  \item Récupérer les solutions dans le récipient \og Tri des déchets Acide et Base \fg{}.
  \item Nettoyer puis ranger le matériel.
\end{itemize}

\begin{tcolorbox}[colback=orange!10!white, colframe=orange!80!black]
\textbf{Appeler le professeur.}
\end{tcolorbox}

\end{document}
